\section*{Pictures of Task Definitions}
\label{appendix:task_def}

\begin{figure}[h!]
    \centering

    % Erste Zeile
    \begin{subfigure}{0.45\textwidth}
        \centering
        \includegraphics[width=0.8\linewidth, angle=90]{task1.png}
        \caption{\textbf{Semicircle:} right means the inner track and left means the outer loop.}
        \label{fig:semicircle}
    \end{subfigure}
    \hfill
    \begin{subfigure}{0.45\textwidth}
        \centering
        \includegraphics[width=0.8\linewidth, angle=90]{task2.jpeg}
        \caption{\textbf{Long corner}: right means the inner track and left means the outer loop.}
        \label{fig:longcorner}
    \end{subfigure}

    \vspace{0.5cm}

    % Zweite Zeile
    \begin{subfigure}{0.45\textwidth}
        \centering
        \includegraphics[width=0.8\linewidth]{4-way.jpg}
        \caption{\textbf{4-way:} use this 4-way intersection for staight, left, right testing.}
        \label{fig:4way}
    \end{subfigure}
    \hfill
    \begin{subfigure}{0.45\textwidth}
        \centering
        \includegraphics[width=0.8\linewidth, angle=-90]{3-way.jpg}
        \caption{\textbf{3-way:} use this 3-way intersection for staight, left, right testing.}
        \label{fig:3way}
    \end{subfigure}

    \caption{Tasks from Experiment section, to test the intents isolated.}
    \label{fig:tasks}
\end{figure}

\begin{figure}[h!]
    \centering
    \includegraphics[width=0.8\linewidth]{loop_task_definition.jpg}
    \caption{\textbf{Long run:} continously driving along this predefined loop.}
    \label{fig:8loop}
\end{figure}

\section*{Network architures}
\label{appendix:architectures}

\begin{figure}[h!]
    \centering
    \includegraphics[width=0.8\linewidth]{pilotnet_4_heads/pilotnet_gated.png}
    \caption{Architecture with multible heads}
    \label{fig:multibleHeads}
\end{figure}

\begin{figure}[h!]
    \centering
    \includegraphics[width=0.8\linewidth]{pilotnet_film/pilotnet_film.png}
    \caption{FiLM architecture}
    \label{fig:FiLM}
\end{figure}


\section*{Other approaches}
\label{subsec:other}

\textcolor{red}{move this to appendix...} \\
In this section, we'd like to briefly describe some other architectures that we tried but that either failed or did not produce the desired results. \\
\textbf{Naive approach.} First we tried to pass the current RGB frame $f_t$ directly a neural network (CNN with fully-connected layers afterwards). 
The intent was fed into the linear layer as a one-hot encoded four dimensional vector. 
We observed, that the model drove forward, but was barely able to react on changes to the lane. This approach likely failed, because the model used the whole image as an input and therefore recieved too much uneccecary information. 
We also assume, that the conditioning method was too weak to have an influence. \\
\textbf{Classification.} Because we used discrete control signals (as described in section \ref{sec:datacollection}) during data collection, we also tried if the model can accurately classify when to do which action.
The model therefore had 3 output neurons, referring to the three actions: 'w' = straight, 'w' + 'a' = left turn, 'w' + 'd' right turn. 
We tried this in combination with the multible head architecture from section \ref{sec:MH}. 
We observed, that the model is able to control the Duckiebot, but also shows a highly non-smooth and zigzagy driving style. (\textcolor{red}{link to a video?})
We therefore decided not to pursue this approach any further.

\section*{VLM results}
\label{sec:vlmresults}

\begin{table}[t]
\centering
\caption{F1 scores for the evaluated approaches.}
\label{tab:classification_results}
\small
\begin{tabular}{lccccc}
\toprule
Metric & Zero-shot & Few-shot & SmolVLM2 & SmolVLM2 & CNN \\
 & Photo & Photo & Mask & Photo & Mask \\
\midrule
Junction        & 0.681 & 0.719 & \textbf{1.000} & 0.970 & 0.836 \\
Left Open       & 0.720 & 0.417 & 0.923 & \textbf{1.000} & 0.720 \\
Straight Open   & 0.222 & 0.000 & \textbf{1.000} & \textbf{1.000} & 0.933 \\
Right Open      & 0.250 & 0.000 & \textbf{1.000} & \textbf{1.000} & 0.609 \\
Exit Cleared    & 0.000 & 0.000 & 0.945 & \textbf{0.963} & 0.612 \\
\midrule
Macro F1        & 0.375 & 0.227 & 0.974 & \textbf{0.987} & 0.742 \\
\bottomrule
\end{tabular}
\end{table}

\newpage
\section*{Contributions}
\textbf{Elia}:
\\
\textbf{Eric}: Collected the dataset, trained the YOLO segmentation model and implemented the FiLM approach.
\\
\textbf{Felix}: Collected the dataset, implemented the training pipeline and trained the multiple head models and the other approaches. 
\\
\textbf{Pankaj}:

All team members contributed to the report and performed the experiments.

\label{appendix:contributions}